\section{Metodologia}
\label{sec:metodologia}

\subsection{Classificação da pesquisa}
\label{subsec:classificacao}

Esta pesquisa classifica-se como \emph{Design Science Research}, paradigma em
que o conhecimento é produzido pela construção e avaliação de um artefato
destinado a resolver um problema identificado em contexto real
\cite{hevner2004design}. A escolha decorre da natureza do trabalho: o problema
foi levantado em campo, a solução é um artefato de software, e a avaliação
acontece contra o processo que o originou.

O percurso segue as seis atividades do modelo de processo de
\citaautor{peffers2007dsrm}, com entrada centrada no problema --- a origem
declarada quando a pesquisa nasce da observação de uma dificuldade prática, e
não de uma oportunidade tecnológica. As atividades são identificar o problema,
definir objetivos para a solução, projetar e desenvolver, demonstrar, avaliar e
comunicar. No momento da redação deste projeto, a primeira está concluída, a
segunda parcialmente, a terceira em curso, e as três últimas planejadas para a
etapa seguinte.

A condução observa os três ciclos descritos por \citaautor{hevner2007ciclos}. O
ciclo de relevância liga o trabalho ao ambiente de aplicação, e aqui é a
prefeitura estudada: dela vieram o problema, os requisitos e os critérios de
aceitação, e é para ela que o artefato volta para ser avaliado. O ciclo de rigor
liga o trabalho à base de conhecimento, da qual se extraem a literatura de
recuperação de informação e de avaliação, e a que se devolve a contribuição
declarada adiante. O ciclo de projeto alterna construção e avaliação internas,
e é onde a bateria de modelos de representação acontece.

\subsubsection{Posicionamento da contribuição}
\label{subsec:posicionamento}

\citaautor{gregor2013posicionamento} propõem posicionar a contribuição de uma
pesquisa de projeto em duas dimensões: a maturidade do domínio da aplicação e a
maturidade das soluções disponíveis. Este trabalho situa-se no quadrante que os
autores denominam \emph{exaptation}: estender solução conhecida a um problema
novo.

A recuperação por proximidade de significado é tecnologia madura, com literatura
consolidada e aplicações em produção. O domínio, porém, é pouco explorado: o
levantamento registrado na seção~\ref{subsec:trabalhos-relacionados} não
encontrou trabalho que apoie a decisão, item a item, de quem compõe uma pesquisa
de preços, e nenhuma das ferramentas em operação examinadas busca o item por
significado.

Os mesmos autores advertem que, nesse quadrante, não basta transpor a técnica:
é preciso demonstrar que o novo campo apresenta dificuldades ausentes naquele
onde ela já era aplicada. Duas o caracterizam. A primeira é que a busca é sempre
submetida a filtro estruturado obrigatório --- categoria, abrangência geográfica
e janela de doze meses, por força de norma ---, e não é o caso irrestrito em que
a literatura costuma medir. A segunda é a natureza do texto: descrições de item
de compra pública têm mediana de seis palavras e seguem gramática própria, o que
as distancia das sentenças sobre as quais os modelos de representação são
habitualmente avaliados.

\subsubsection{Contribuição declarada}
\label{subsec:contribuicao}

A tipologia de \citaautor{gregor2013posicionamento} distingue três níveis de
contribuição, do mais situado ao mais abstrato: a implementação do artefato, o
conhecimento na forma de princípios operacionais, e a teoria de projeto
plenamente desenvolvida. Este trabalho pretende contribuir nos dois primeiros, e
não no terceiro.

No primeiro nível está o sistema, instanciação do artefato no município
estudado. No segundo, e é esta a contribuição que o trabalho reivindica como
principal, está o \textbf{protocolo de avaliação que corrige o viés de origem do
conjunto de referência}.

O problema que esse protocolo trata é o seguinte. O conjunto de referência deste
trabalho é extraído de pesquisas de preço já realizadas, e os preços que o
compõem foram escolhidos pela servidora dentro de uma ferramenta que busca por
palavra. Os itens tidos como relevantes são, por construção, os que uma busca
lexical alcança --- e um método semântico que encontre item comparável descrito
com outras palavras é contado como erro. A literatura de avaliação documenta o
efeito: conjuntos montados apenas com sistemas de uma família avaliam mal os de
outra \pendente{citar Yilmaz et al. 2020 após conferir a entrada}.

A correção proposta reúne os primeiros resultados de todos os métodos
comparados, incluindo a busca lexical, embaralha-os e submete o conjunto a
julgamento sem identificação de origem. A assimetria que ela desfaz está
descrita na seção~\ref{subsec:limites-avaliacao}.

\pendente{Esta declaração de contribuição precisa ser confrontada com a
literatura de avaliação antes de ser afirmada como nova. Há trabalho publicado
tratando de anotação parcial e de viés de coleta, e a redação correta é
posicionar a contribuição como aplicação e adaptação desse conhecimento ao caso
de descrições curtas sob filtro restritivo --- não como invenção.}

\subsubsection{Classe de problemas}
\label{subsec:classe-de-problemas}

O artefato é construído para um município, mas a classe de problemas a que ele
responde é mais ampla, e delimitá-la é o que permite dizer a quem mais o
resultado serve.

A classe é: \textbf{triagem de candidatos a referência, em acervo público de
descrições curtas e heterogêneas, sob filtro estruturado obrigatório, com
decisão final humana e exigência de registro da procedência.} Pertencem a ela os
casos em que a busca precisa recuperar por significado mas obedecer a restrições
formais, o julgamento não pode ser automatizado por imposição normativa, e cada
escolha precisa ser justificável perante terceiro.

Pesquisa de preços em contratação pública é um caso dessa classe. Outros,
não examinados aqui, provavelmente também o são.

\pendente{Confirmar, na literatura, se a classe de problemas enunciada já tem
nome estabelecido. Nomear classe que já existe com outro rótulo é erro que a
revisão apanha.}

\subsection{Procedimento de levantamento}
\label{subsec:procedimento-levantamento}

O levantamento combinou três fontes de evidência, com pesos diferentes.

A primeira foi uma entrevista estruturada com a servidora responsável pela fase
preparatória no setor de licitações do município estudado, seguida de uma sessão
de validação na qual a especificação já redigida foi lida bloco a bloco para que
ela a corrigisse. A segunda foi a análise de documentos de processos encerrados
cedidos pelo setor: documentos de formalização de demanda, estudos técnicos
preliminares, termos de referência e relatórios de cotação. A terceira foi a
medição desses documentos, descrita adiante.

A distinção entre as três importa e é mantida ao longo do texto. O que a
servidora declarou aparece como declaração; o que os documentos mostram aparece
como medida; e o que foi inferido de um pelo outro aparece como inferência,
sempre identificada. A sessão de validação existiu precisamente para separar
essas camadas, e derrubou três afirmações que até então figuravam como
diagnóstico.

A mais consequente das três merece registro, porque ilustra o risco do método.
A comparação entre as datas de assinatura sugeria que o estudo técnico
preliminar era assinado depois da cotação e ainda assim se referia a ela como
evento futuro, o que foi lido como prova documental de inversão da ordem legal.
A leitura estava errada: existe uma minuta assinada antes, e é dela que vem o
tempo verbal futuro. O desvio real é outro e menor, e está descrito na
subseção~\ref{subsec:processo-hoje}. Anomalia de documento só vira evidência
depois de verificado quantas versões daquele documento existem.

\pendente{Detalhar, quando a avaliação for desenhada por completo, o tamanho e o
critério de seleção do conjunto de documentos analisados, e explicitar que ele é
de conveniência.}

\subsection{Tratamento dos dados de domínio}

Documentos reais de processos licitatórios não são versionados nem reproduzidos
no texto. Cada documento recebido é mantido fora do repositório e dele se deriva
uma versão sintética, com a mesma estrutura e nenhum dado verdadeiro, que é a
única utilizada em testes, exemplos e ilustrações deste trabalho.

\pendente{Detalhar o procedimento de derivação e discutir suas limitações para
a avaliação da solução.}

\subsection{Desenho da solução}

O desenho que resultou do levantamento --- requisitos, regras de negócio e
arquitetura --- é apresentado na seção~\ref{sec:requisitos-arquitetura}, e não
aqui. A separação tem motivo: ele é resultado do percurso metodológico e não
etapa dele, e tratá-lo como procedimento esconderia que boa parte das escolhas
foi determinada por evidência de campo, não por preferência técnica.

Duas escolhas, porém, pertencem à metodologia, porque dizem respeito à natureza
da evidência que as sustenta.

\subsubsection{Recuperação de informação sobre a legislação e os modelos}

\pendente{Recuperação aumentada por geração aplicada à legislação e aos modelos
de documento já aprovados. Descrever a estratégia de fragmentação, o modelo de
representação escolhido e o critério de recuperação. Há material reaproveitável
dos fichamentos da Iniciação Científica, mas cada citação precisa ser
reconferida contra o PDF.}

\pendente{Posicionar contra o contra-argumento levantado no levantamento de
literatura: há trabalho publicado sustentando que recuperação plana sobre
documento não estruturado limita raciocínio de múltiplos saltos quando os dados
são estruturados. Responder explicitamente, não por omissão.}

\subsubsection{Natureza da evidência que sustenta a escolha da persistência}

A separação entre banco relacional e índice semântico, descrita na
subseção~\ref{subsec:arquitetura}, apoia-se em medição realizada pelo autor em
iniciação científica anterior a este trabalho, e não em literatura revisada por
pares. O texto registra isso onde a medição é invocada, e a distinção é mantida
deliberadamente: medição própria é evidência, com força menor que a de resultado
publicado e submetido a revisão.

\pendente{Descrever o protocolo daquela medição --- conjunto de dados,
seletividade do filtro, métrica de recuperação e número de repetições --- em
detalhe suficiente para que o leitor julgue se ela transfere para o caso deste
trabalho. Critério acrescentado pela orientação: verificar contra a natureza dos
dados a tratar, não só o desempenho medido.}

\subsection{Processo de desenvolvimento}

\pendente{Processo de desenvolvimento, com testes automatizados como critério
de conclusão de cada incremento.}

\subsection{Avaliação}

\pendente{Definir como a solução será avaliada e contra o quê. Sem uma linha de
base medida antes da intervenção, não há como afirmar redução de esforço ---
planejar a coleta dessa linha de base ainda em TCC I.}

\pendente{A etapa de verificação pelo órgão de controle ficou fora do escopo
obrigatório por decisão da orientação em 19/08/2026: o tempo de resposta do
órgão não é previsível e não cabe em planejamento. Reaproveitável no fim, se
houver tempo.}
